% AUTORIDAD DE FUENTE: Exposé IX §9.3, líneas francesas 2279--2430;
% índices de escaneo 432--438, folios 421--427. Las ocho operaciones
% materiales indicadas por las capas D/C se conservan mediante \sisi.

\medskip
\noindent
9.3. \underline{Extensiones mixtas.}

Sea $\underline C$ una categoría abeliana, y sean
\begin{equation}\tag{9.3.1}
 P,\ R,\ Q
\end{equation}
tres objetos de $\underline C$. Nos interesa clasificar los objetos $X$
de $\underline C$ dotados de una filtración de tres escalones
\begin{equation}\tag{*}
 X^0=X\supset X^1\supset X^2\supset X^3=0
\end{equation}
y de isomorfismos dados
\begin{equation*}
 X^0/X^1\xrightarrow{\ \sim\ }\sisi{P}{Q},
 \qquad X^1/X^2\xrightarrow{\ \sim\ }R,
 \qquad X^2/X^3=X^2\xrightarrow{\ \sim\ }\sisi{Q}{P}.
\end{equation*}
Poniendo
\begin{equation*}
 E=X^0/X^2,
 \qquad F=X^1/X^3=X^1,
\end{equation*}
vemos que $X$ da lugar a dos extensiones ordinarias
\begin{equation}\tag{9.3.2}
 0\longrightarrow P\longrightarrow F\longrightarrow R\longrightarrow0,
 \qquad
 0\longrightarrow R\longrightarrow E\longrightarrow Q\longrightarrow0,
\end{equation}
de $R$ por $P$ y de $Q$ por $R$, respectivamente. Nuestro punto de
vista consistirá en partir, además de los objetos (9.3.1), de las
estructuras de extensión (9.3.2), que también se suponen dadas, y
considerar las estructuras filtradas correspondientes sobre $X$.

Más precisamente, una \underline{extensión mixta} de la extensión $E$
por la extensión $F$ será un objeto $X$ de $\underline C$, dotado de
una filtración de tres escalones (*) y de isomorfismos dados
\begin{equation}\tag{9.3.3}
 E\simeq X^0/X^2,
 \qquad F\simeq X^1/X^3=X^1,
\end{equation}
compatibles con las filtraciones: llevan $X^1/X^2$ a
$\operatorname{Im}(R\longrightarrow E)$ y $X^2$ a
$\operatorname{Im}(P\longrightarrow F)$, respectivamente, y se exige
que el isomorfismo $R\simeq F/P$ —y, por tanto,
$R\simeq X^1/X^2$— deducido de (9.3.3) sea el deducido de la primera
sucesión exacta de (9.3.2).

%%%%%%%%%%  scan idx 433  |  folio 422  %%%%%%%%%%

Para $P,Q,R$ fijos y extensiones $E,F$ fijas, las extensiones mixtas de
$E$ por $F$ forman una categoría de manera evidente; de hecho, es un
grupoide, pues toda flecha es invertible. El grupo de automorfismos de
una extensión mixta $X$ de $E$ por $F$ es canónicamente isomorfo a
$\Hom(Q,P)$ mediante
\begin{equation}\tag{9.3.4}
 \Hom(Q,P)
 \xrightarrow{\ \sim\ }
 \operatorname{Aut}_{\mathrm{extpan}}(X),
 \qquad
 u\longmapsto\id_X+\bar u,
\end{equation}
donde $\bar u$ es el endomorfismo del objeto subyacente a $X$ obtenido
componiendo
$X\longrightarrow Q\xrightarrow{u}P\longrightarrow X$; las flechas
exteriores son, respectivamente, el epimorfismo canónico y el
monomorfismo canónico provenientes de la estructura de extensión mixta
dada sobre $X$.

Se puede precisar la estructura de la categoría
$\operatorname{EXTPAN}(E,F)$ de extensiones mixtas de $E$ por $F$
considerándola como un ``torsor'' bajo la ``categoría-grupo'' (VII 2.5)
$\operatorname{EXT}(Q,P)$ de extensiones de $Q$ por $P$, mediante un
funtor natural análogo a la composición de Brauer de extensiones
ordinarias:
\begin{equation}\tag{9.3.5}
\begin{aligned}
 \omega:
 \operatorname{EXT}(Q,P)\times\operatorname{EXTPAN}(E,F)
 &\longrightarrow\operatorname{EXTPAN}(E,F),\\
 \omega(G,X)&=G\wedge X.
\end{aligned}
\end{equation}
Para todo objeto $X$ de $\operatorname{EXTPAN}(E,F)$, el funtor
parcial $G\longmapsto G\wedge X$ es una equivalencia de
$\operatorname{EXT}(Q,P)$ con $\operatorname{EXTPAN}(E,F)$.

%%%%%%%%%%  scan idx 434  |  folio 423  %%%%%%%%%%

Daremos solo una breve descripción de $\omega$. Consideremos una
extensión mixta $X$ de $E$ por $F$ como una extensión ordinaria, también
denotada $X$, de $Q$ por $F$. Para toda extensión $G$ de $Q$ por $P$,
sea $\bar G$ la extensión de $Q$ por $F$ inducida por la inclusión
$P\longrightarrow F$. Pongamos
\begin{equation}\tag{9.3.6}
 G\wedge X=\omega(G,X)
 \overset{\mathrm{def}}{=}\bar G\wedge X,
\end{equation}
donde $\wedge$ denota la composición de Brauer en la categoría de
extensiones de $Q$ por $F$. La filtración de tres escalones del miembro
derecho $X'$ de (9.3.6) es clara, al igual que el segundo isomorfismo de
tipo (9.3.3) asociado a $X'$.

Para definir el primer isomorfismo de tipo (9.3.3) para $X'$, obsérvese
que el funtor
\begin{equation*}
 Y\longmapsto\widetilde Y,
 \qquad
 \operatorname{EXT}(Q,F)\longrightarrow\operatorname{EXT}(Q,R),
\end{equation*}
inducido por el morfismo dado $F\longrightarrow R$ de (9.3.2) es
compatible con la composición de Brauer. Por tanto, envía $X'$ a una
extensión $\widetilde X'$ canónicamente isomorfa a
$\widetilde{\bar G}\wedge\widetilde X$. El primer isomorfismo (9.3.3)
da un isomorfismo
$\widetilde X\simeq\sisi{F}{E}$, mientras que
$\widetilde{\bar G}$ —la extensión inducida por $G$ mediante el
compuesto $P\longrightarrow F\longrightarrow R$— es canónicamente
isomorfa a la extensión trivial, pues ese compuesto es nulo. En
consecuencia, existe un isomorfismo canónico
$\widetilde X'\simeq\sisi{F}{E}$, que especifica sobre el miembro
derecho de (9.3.6) su estructura de extensión mixta de $E$ por $F$.

%%%%%%%%%%  scan idx 435  |  folio 424  %%%%%%%%%%

Dejamos al lector la definición de $\omega$ sobre las flechas de
$\operatorname{EXT}(Q,P)\times\operatorname{EXTPAN}(E,F)$, la
verificación de que (9.3.5) posee la propiedad enunciada —para toda
extensión mixta \underline{fija} $X$, el funtor
$G\longmapsto G\wedge X$ es una equivalencia de
$\operatorname{EXT}(Q,P)$ con $\operatorname{EXTPAN}(E,F)$—, y la
demostración de la fórmula de asociatividad
\begin{equation}\tag{9.3.7}
 G_1\wedge(G_2\wedge X)
 \simeq(G_1\wedge G_2)\wedge X,
\end{equation}
donde $G_1\wedge G_2$ denota la composición de Brauer en
$\operatorname{EXT}(Q,P)$.

La discusión precedente da inmediatamente (a) y (b) del resultado
siguiente.

\medskip
\noindent
\underline{Proposición 9.3.8. Dadas dos extensiones} $E$ \underline{y} $F$ \underline{como en
(9.3.2) en la categoría abeliana} $\mathbf C$, \underline{consideremos la
categoría} $\operatorname{EXTPAN}(E,F)$ \underline{de extensiones mixtas de} $E$ \underline{por}
$F$. \underline{Entonces:}

\smallskip
\noindent
a) \underline{$\operatorname{EXTPAN}(E,F)$ es un grupoide. Para cada
uno de sus objetos, el grupo de automorfismos es canónicamente
isomorfo a}
$\Ext^0(Q,P)=\Hom(Q,P)$.

\smallskip
\noindent
b) \underline{El conjunto $\operatorname{Extpan}(E,F)$ de clases de
isomorfismo de objetos de $\operatorname{EXTPAN}(E,F)$ es vacío o es,
de manera natural, un torsor bajo el grupo} $\Ext^1(Q,P)$.

\smallskip
\noindent
c) \underline{Se puede definir canónicamente una clase}
\begin{equation}\tag{9.3.9}
 c(E,F)\in\Ext^2(Q,P),
\end{equation}
\underline{cuya anulación es necesaria y suficiente para que}
$\operatorname{Extpan}(E,F)\neq\varnothing$, \underline{es decir, para
que exista una extensión mixta de} $E$ \underline{por} $F$.

%%%%%%%%%%  scan idx 436  |  folio 425  %%%%%%%%%%

Queda por demostrar $\sisi{\mathrm{d})}{\mathrm{c})}$. Una extensión
mixta de $E$ por $F$ puede describirse también como una extensión $X$
de $E$ por $P$ junto con un isomorfismo de extensiones de $R$ por $P$
\begin{equation*}
 X^1\simeq F,
\end{equation*}
donde $X^1$ denota la imagen de la extensión $X$ por el morfismo
$R\longrightarrow E$ de (9.3.2). Por tanto, tal extensión mixta existe
si y solo si hay un elemento de $\Ext^1(E,P)$ cuya imagen en
$\Ext^1(R,P)$ sea la clase $\xi$ de $F$. Consideremos la sucesión
exacta de los $\Ext^i(-,P)$ asociada a la segunda sucesión exacta de
(9.3.2):
\begin{equation*}
 \cdots\longrightarrow
 \Ext^1(E,P)\longrightarrow
 \Ext^1(R,P)
 \xrightarrow{\partial}
 \sisi{\Ext^2(R,Q)}{\Ext^2(Q,P)}\longrightarrow\cdots .
\end{equation*}
Basta definir $c(E,F)$ por
\begin{equation}\tag{9.3.10}
 c(E,F)=\partial\xi,
 \qquad
 \xi=\operatorname{cl}(F)\in\Ext^1(R,P).
\end{equation}

Supongamos ahora dado un funtor exacto
\begin{equation}\tag{9.3.11}
 r:\underline C\longrightarrow\underline C',
\end{equation}
donde $\underline C'$ es otra categoría abeliana. Denotemos por
$P',Q',R',E',F'$ los objetos de $\underline C'$ obtenidos aplicando
$r$ a $P,Q,R,E,F$, de modo que $E'$ sea una extensión de $Q'$ por $R'$
y $F'$ una extensión de $R'$ por $P'$. Salvo isomorfismo canónico, el
diagrama de funtores siguiente es conmutativo:
\begin{equation}\tag{9.3.12}
\begin{tikzcd}[column sep=large,row sep=large]
 \operatorname{EXT}(Q,P)\times\operatorname{EXTPAN}(E,F)
   \arrow[r,"\omega"]\arrow[d]
 &\operatorname{EXTPAN}(E,F)\arrow[d]\\
 \operatorname{EXT}(Q',P')\times\operatorname{EXTPAN}(E',F')
   \arrow[r,"\sisi{\omega}{\omega'}"']
 &\operatorname{EXTPAN}(E',F').
\end{tikzcd}
\end{equation}
Las flechas verticales son los funtores evidentes inducidos por el
funtor exacto $r$, y $\omega'$ es la flecha análoga a $\omega$ de
(9.3.5), sustituyendo $\underline C$ por $\underline C'$. Se obtienen
compatibilidades análogas con $r$ para las estructuras consideradas en
9.3.8; dejamos su formulación explícita al lector.

%%%%%%%%%%  scan idx 437  |  folio 426  %%%%%%%%%%

\begingroup
\sloppy
Utilizaremos la consecuencia siguiente de esta conmutatividad.
Supongamos que $\operatorname{EXTPAN}(E,F)$ no es vacía y elijamos un
objeto $X$ en ella. Mediante $G'\longmapsto G'\wedge X'$, donde
$X'=r(X)$, la categoría $\operatorname{EXTPAN}(E',F')$ es equivalente
a $\operatorname{EXT}(Q',P')$. Así, un objeto $Y'$ de
$\operatorname{EXTPAN}(E',F')$ queda determinado salvo isomorfismo
cuando se conoce la clase de isomorfismo del objeto correspondiente de
$\operatorname{EXT}(Q',P')$, es decir, cuando se conoce el elemento
correspondiente
\par
\endgroup
\begin{equation*}
 c_X(Y')\in\Ext^1(Q',P').
\end{equation*}

Si $X$ se sustituye por otro objeto $X_1=H\wedge X$, donde $H$ es una
extensión de $Q$ por $P$, se obtiene inmediatamente de (9.3.7)
\begin{equation*}
 c_{X_1}(Y')=c_X(\sisi{Y}{Y'})-\operatorname{cl}(H').
\end{equation*}
Por tanto, para $Y'$ fijo y $X$ variable, los elementos $c_X(Y')$
recorren exactamente una clase lateral
\begin{equation}\tag{9.3.13}
 c(Y')\in
 \Ext^1(Q',P')/\operatorname{Im}\Ext^1(Q,P),
\end{equation}
asociada canónicamente a la extensión mixta $Y'$ de $E'$ por $F'$. Por
construcción, su anulación es necesaria y suficiente para que $Y'$ sea
isomorfa a la transformada por $r$ de una extensión mixta $X$ de $E$
por $F$.
